This appendix provides supplementary mathematical background for the key techniques employed in TeachWise.

\section*{Genetic Algorithm Fitness Function}

The fitness of a candidate schedule $S$ is computed as:
\[
f(S) = \sum_{a \in S} \Bigl[ w_{\text{hard}} \cdot \delta_{\text{hard}}(a) + w_{\text{med}} \cdot \delta_{\text{med}}(a) + w_{\text{soft}} \cdot \phi_{\text{soft}}(a) \Bigr]
\]
where $\delta_{\text{hard}}(a)$ and $\delta_{\text{med}}(a)$ are penalty indicator functions for hard and medium constraint violations respectively, and $\phi_{\text{soft}}(a)$ is the soft-constraint quality score for assignment $a$. The penalty weights are set as $w_{\text{hard}} = -50{,}000$, $w_{\text{med}} = -500$, and soft-constraint bonuses include $+500$ per satisfied student time preference and $+100 \times m$ for a pedagogical match score $m \in [0,1]$.

\section*{Prosodic Feature Extraction}

Let $x(t)$ denote a raw audio waveform sampled at 16,kHz. The fundamental frequency (pitch) $F_0(t)$ is estimated frame-by-frame using the YIN algorithm. Pitch variability is captured by the standard deviation $\sigma_{F_0}$ over a 3--4 minute segment. Loudness is computed as the root-mean-square energy:
[
$\text{RMS}(t) = \sqrt{\frac{1}{N}\sum_{n=0}^{N-1} x^2(t+n)}$
]

Additional acoustic features, including jitter (cycle-to-cycle pitch variation) and shimmer (cycle-to-cycle amplitude variation), are extracted using the OpenSMILE \texttt{eGeMAPSv02} feature set. A subset of 18 selected features per segment is used as input to the Random Forest classifiers.

\section*{Random Forest Classification}

Each of the four quality-dimension classifiers (energy, patience, student attention, and lesson pace) is implemented as a Random Forest ensemble consisting of $T = 100$ decision trees. Each tree uses the Gini impurity criterion for splitting and is constrained to a maximum depth of 4 to reduce overfitting and improve generalization. Class imbalance is handled using balanced class weights. At each split, a subset of features is selected using the $\sqrt{d}$ heuristic. Final predictions are obtained through majority voting across all trees:
\[
\hat{y} = \arg\max_{c \in \{1,2,3,4,5\}}
\sum_{t=1}^{T} \mathbb{1}\big(h_t(\mathbf{x}) = c\big)
\]

where $h_t$ represents the prediction of the $t$-th tree and $c \in {1,2,3,4,5}$ is the Likert scale label.


\section*{Curriculum Fidelity Assessment}

The curriculum fidelity score is obtained using a large language model (LLaMA3) via prompt-based evaluation. The lesson transcript and corresponding curriculum content are provided directly as input to the model, which generates a structured rubric score (1--5) based on alignment with curriculum objectives and inclusion of relevant material.

This approach does not employ a retrieval-augmented generation (RAG) pipeline; no embedding-based similarity or document retrieval is performed. Instead, the evaluation is conducted using zero-shot or few-shot prompting.

\section*{Hybrid Teaching Score}

The final composite teaching effectiveness score for a session is an unweighted average of five dimensions:
\[
\text{Score} = \frac{E + P + A + L + C}{5}
\]
where $E$ is Energy, $P$ is Patience, $A$ is Student Attention, $L$ is Lesson Pace (each from the Random Forest classifiers, scaled to 1--5), and $C$ is the RAG-derived Curriculum Fidelity score (1--5).